
This work establishes construct validation as a prerequisite for proxy-based optimization in code generation. We present a four-stage validation pipeline with Stage 1 (measurement reliability) empirically validated via proof-of-concept. The framework tests proxies through coefficient of variation (CV $\leq 5$\%), effect size (Cohen's d $\geq 0.8)$, and cross-platform stability (Spearman $\rho \geq 0.8)$ before reinforcement learning optimization.

Proof-of-concept validation demonstrates that structural similarity (CodeBLEU) achieves exceptional reliability (CV=1.39\%, Cohen's d=4.51, Spearman $\rho$=0.949), passing all three criteria with substantial margins. Runtime efficiency measurements using a wall-clock noise model marginally exceed the CV threshold (6.22\% vs 5.0\%), with literature evidence suggesting hardware instruction counting would achieve lower variance, though empirical confirmation is needed.

The framework employs a scoped gate design where partial validation ($\geq 1$ proxy passing all criteria) constitutes progression. This prevents all-or-nothing failure and recognizes that identifying which proxies are reliable---and which require specialized instrumentation---is scientifically valuable.

Our contribution shifts code generation evaluation from correlation-based heuristics to rigorous reliability testing. Different quality dimensions (structural similarity, runtime efficiency, style conformity) have distinct measurement reliability profiles requiring independent validation. Before investing thousands of GPU hours in multi-objective reinforcement learning, we must validate that quality proxies are measurable---a prerequisite this work establishes.

Future work includes: (1) real infrastructure validation (CodeLlama-7B + HumanEval + hardware counters), (2) Stage 2 conditional independence testing via hierarchical regression, (3) cross-repository generalization validation, and (4) full multi-objective RL training with validated proxies. Whether validated proxies yield Pareto improvements or fail conditional independence, the outcome advances understanding by providing the prerequisite framework the field lacked.
